% Generated draft: requires mathematical and source review.
\documentclass[11pt]{article}
\usepackage[margin=25mm]{geometry}
\usepackage{fontspec}
\setmainfont{Latin Modern Roman}
\usepackage{amsmath,amssymb,mathtools,unicode-math}
\setmathfont{Latin Modern Math}
\usepackage{enumitem,needspace,xurl,hyperref}
\setlength{\parindent}{0pt}
\setlength{\parskip}{5pt}
\setlength{\emergencystretch}{3em}
\newcommand{\N}{\mathbb{N}}
\newcommand{\Z}{\mathbb{Z}}
\newcommand{\Q}{\mathbb{Q}}
\newcommand{\R}{\mathbb{R}}
\newcommand{\C}{\mathbb{C}}
\newcommand{\F}{\mathbb{F}}
\newcommand{\A}{\mathbb{A}}
\newcommand{\PP}{\mathbb{P}}
\newcommand{\E}{\mathbb{E}}
\newcommand{\eps}{\varepsilon}
\newcommand{\dd}{\,\mathrm{d}}
\begin{document}
\section*{UnsolvedMath ID: 30005990}
\textbf{Source problem code:} catalogue entry
\subsection{Definitions and mathematical statement}
Let \(d\in\N\), \(d>2\). The supplied record concerns a spectral optimal partition of a domain \(D\subset\R^d\), its associated eigenfunctions \(u_1,\ldots,u_N\), and a free interface \(\mathcal F\). A boundary free-interface point is intended to be a point where the free interface reaches \(\partial D\); the precise definition of this boundary contact set is not supplied. Write \(\gamma(x_0)\) for the limiting frequency at such a point, whenever the relevant frequency function and its limit are defined. The original excerpt contains a corrupted integral formula, so neither the exact frequency function nor its normalization can be recovered reliably from this record.

For a fully specified admissible class of such partition problems in dimension \(d\), let
\[
\Gamma_d=\{\gamma(x_0):x_0\text{ is an admissible boundary free-interface point}\}.
\]
Interpreting \(\delta_d\) as the first possible frequency level strictly above \(2\), define conditionally
\[
\delta_d:=\inf\bigl(\Gamma_d\cap(2,\infty)\bigr),
\qquad \inf\varnothing:=+\infty.
\]
This interpretation makes the conjectural value in the supplied statement meaningful; the record does not provide an uncorrupted definition of \(\delta_d\). In particular, this quantity is a frequency threshold, whereas the corresponding numerical separation from \(2\) would be \(\delta_d-2\).

\textbf{Conjectural target, conditional on the missing admissibility and frequency definitions.} For every integer \(d>2\), determine \(\delta_d\), and prove or disprove that
\[
\Gamma_d\cap(2,5/2)=\varnothing,
\qquad \frac52\in\Gamma_d,
\qquad\text{hence}\qquad \delta_d=\frac52.
\]
The supplied conjecture further specifies that the value \(5/2\) is attained by three partition domains forming a triple junction of \(3/2\)-points arranged along a line that reaches \(\partial D\) orthogonally. The record does not define a \(3/2\)-point or specify whether this geometry describes the original partition, a tangent configuration, or a homogeneous blow-up. Those alternatives must be resolved before the attainment clause becomes a precise geometric assertion.

The spectral objective, admissible partition class, number of phases, boundary conditions, domain hypotheses, frequency definition, and boundary contact notion are missing. Consequently, the displayed formulation records a conditional target rather than a complete unconditional conjecture. The excerpt emphasizes technical and topological boundary difficulties: frequency monotonicity is problematic even for \(C^{1,\alpha}\) boundaries; defining the optimal partition on the boundary requires analysis of eigenfunctions near \(\partial D\); and regular free-boundary smoothness is described as relying on blow-up decay obtained through epiperimetric inequalities, rather than an implicit-function-theorem argument. These comments do not themselves specify the hypotheses of the conjecture.
\subsection{Short English statement}
For spectral optimal partitions in dimension \(d>2\), determine the sharp boundary frequency threshold above \(2\). The supplied conjecture predicts the threshold \(5/2\), attained by a three-domain junction whose \(3/2\)-points lie along a line meeting the boundary orthogonally. The admissible problems, frequency normalization, and precise junction geometry remain unspecified, so this is a conditional statement-only draft requiring mathematical review.
\subsection{Sources}
\begin{itemize}
\item Supplied local UnsolvedMath registry record, permanent ID \(30005990\), title \emph{Boundary Frequency Gap for Optimal Partitions}, exact source problem code \texttt{catalogue entry}; collection: \emph{Oberwolfach Reports Open Problems}. The record supplies proposed year \(2024\), but no named proposer.
\item Supplied report citation: \emph{Calculus of Variations (2025)}, Oberwolfach Reports, DOI \texttt{10.4171/owr/2024/37}, pp.\ 2113--2178. Source URL and supplied literature-source link, labeled ``Oberwolfach source report'': \url{https://doi.org/10.4171/owr/2024/37}. This report is supplied as primary provenance; its contents have not been independently consulted here. The citation's \(2025\) label and its \(2024\) DOI component are retained as supplied.
\item Literature assessment dated exactly \texttt{2026-08-22}: the record labels the problem open and reports that only a positive dimension-dependent boundary frequency gap is known, while the conjectural threshold \(5/2\) remains unproved and is listed as open in the cited report. This is supplied dated curation, not an independently verified mathematical or literature assessment.
\item The supplied registry labels the statement \texttt{corrected\_verified} and says that it was reconstructed and checked against the official Oberwolfach report. This is a source-provided verification claim, not verification performed for this draft. The damaged excerpt refers to reference [8], but its bibliographic entry is absent.
\end{itemize}
\end{document}
